% ===== PRÉAMBULE CANONIQUE — à coller VERBATIM en tête de CHAQUE .tex =====
\documentclass[11pt,a4paper]{article}
\usepackage[utf8]{inputenc}\usepackage[T1]{fontenc}\usepackage{lmodern}
\usepackage[french]{babel}
\usepackage{amsmath,amssymb,mathtools}\usepackage{icomma}
\usepackage{siunitx}\sisetup{locale=FR}  % \qty{}{}, \si{}, \ang{}
\usepackage{xcolor}\usepackage{colortbl}
\usepackage{tikz}\usepackage{pgfplots}\pgfplotsset{compat=1.18}
\usepackage[siunitx,europeanresistors]{circuitikz} % circuits (élec.)
\usepackage[most]{tcolorbox}\usepackage{enumitem}\usepackage{array,booktabs}
\usepackage[a4paper,margin=2.2cm]{geometry}
\usetikzlibrary{arrows.meta,calc,angles,quotes,positioning,patterns,%
  backgrounds,shapes.geometric,decorations.pathmorphing,decorations.markings,intersections}
\definecolor{primary}{RGB}{222,49,99}\definecolor{primarydark}{RGB}{150,22,55}
\definecolor{defcol}{RGB}{0,114,178}\definecolor{methcol}{RGB}{52,92,148}
\definecolor{excol}{RGB}{0,135,90}\definecolor{attncol}{RGB}{200,40,55}
\tcbset{boxbase/.style={enhanced,breakable,boxrule=0.9pt,arc=2.5pt,
  left=3mm,right=3mm,top=2.3mm,bottom=2.3mm,fonttitle=\bfseries,coltitle=white}}
% --- boîtes TUTORIEL (noms EXACTS, ne pas renommer) ---
\newtcolorbox{cle}[1][]{boxbase,colback=defcol!6,colframe=defcol,
  title={\textsc{À retenir}\ifx\relax#1\relax\relax\else\textmd{ : \itshape #1}\fi}}
\newtcolorbox{methode}[1][]{boxbase,colback=methcol!7,colframe=methcol,sharp corners,
  title={\textsc{Méthode}\ifx\relax#1\relax\relax\else\textmd{ --- \itshape #1}\fi}}
\newtcolorbox{exemplebox}[1][]{boxbase,colback=excol!5,colframe=excol,
  title={\textsc{Exemple}\ifx\relax#1\relax\relax\else\textmd{ : \itshape #1}\fi}}
\newtcolorbox{piege}[1][]{boxbase,colback=attncol!6,colframe=attncol,
  title={\textsc{Piège}\ifx\relax#1\relax\relax\else\textmd{ : \itshape #1}\fi}}
\newtcolorbox{remarque}[1][]{boxbase,colback=black!4,colframe=black!45,
  title={\textsc{Remarque}\ifx\relax#1\relax\relax\else\textmd{ : \itshape #1}\fi}}
% --- boîtes EXERCICE / CORRIGÉ (noms EXACTS, auto-comptées) ---
\newtcolorbox[auto counter]{exo}[1][]{boxbase,colback=primary!4,colframe=primary,
  title={\textsc{Exercice}~\thetcbcounter\ifx\relax#1\relax\relax\else\textmd{ --- \itshape #1}\fi}}
\newtcolorbox[auto counter]{corrige}[1][]{boxbase,colback=excol!4,colframe=excol,
  title={\textsc{Corrigé}~\thetcbcounter\ifx\relax#1\relax\relax\else\textmd{ --- \itshape #1}\fi}}
\newcommand{\vv}[1]{\vec{#1}}\newcommand{\ds}{\displaystyle}
\newcommand{\R}{\mathbb{R}}\newcommand{\N}{\mathbb{N}}\newcommand{\Z}{\mathbb{Z}}
% (un \usepackage{fancyhdr}+en-tête est possible ; il est ignoré côté web)
% =========================================================================
\begin{document}
\sisetup{per-mode=symbol}

\begin{corrige}[vrai ou faux : combien de propositions correctes ?]
\begin{enumerate}[label=\alph*)]
  \item \textbf{Faux}. On ne crée jamais de charge : on transfère des électrons (conservation).
  \item \textbf{Vrai}. $q = N e$ avec $N\in\Z$ : la charge est quantifiée.
  \item \textbf{Faux}. Un métal nu ne se charge pas durablement par frottement (charges libres
        aussitôt redistribuées / écoulées).
  \item \textbf{Vrai}. Un corps chargé polarise un neutre léger par influence, puis l'attire.
  \item \textbf{Vrai}. Le nombre de neutrons est indépendant de la charge ; un corps neutre a
        seulement $N_e = N_p$, sans contrainte sur les neutrons.
  \item \textbf{Faux}. Dans un isolant les charges sont \emph{liées} : elles ne migrent pas pour
        combler un défaut (c'est justement ce qui distingue un isolant d'un conducteur).
\end{enumerate}
\textbf{Bilan : 3 propositions correctes} (b, d, e).
\end{corrige}

\begin{corrige}[l'électroscope sous influence]
\begin{enumerate}[label=\alph*)]
  \item La baguette \emph{positive} attire les électrons libres \textbf{vers la boule}. La boule,
        enrichie en électrons, devient \emph{moins positive} : sa charge \textbf{diminue}. Les
        électrons ayant quitté la région des feuilles, celles-ci deviennent \emph{plus positives} :
        elles se repoussent davantage et l'écartement \textbf{augmente}.
  \item Une baguette \emph{négative} \emph{repousse} les électrons vers le bas, vers les feuilles.
        La boule devient \emph{plus positive} (sa charge augmente), tandis que les feuilles,
        partiellement neutralisées par ces électrons, deviennent \emph{moins positives} :
        l'écartement \textbf{diminue}.
\end{enumerate}
\end{corrige}

\begin{corrige}[deux barres sous influence]
\begin{enumerate}[label=\alph*)]
  \item L'isolant positif \emph{attire} les électrons libres vers l'extrémité proche : ils
        s'accumulent dans $A$. La barre $A$ (proche) devient donc \textbf{négative}, et $B$
        (éloignée), appauvrie en électrons, devient \textbf{positive}. En séparant les barres
        \emph{avant} d'ôter le bâton, ces charges restent piégées : $q_A<0$, $q_B>0$.
  \item $q_A + q_B = 0$ : aucun électron n'a été apporté ni retiré à l'ensemble (pas de contact
        avec l'isolant), c'est le \textbf{principe de conservation de la charge}. Les électrons ont
        seulement migré de $B$ vers $A$.
\end{enumerate}
\end{corrige}

\begin{corrige}[deux pendules électrostatiques]
\begin{enumerate}[label=\alph*)]
  \item Les deux forces sont les deux faces d'une \emph{même} interaction de Coulomb :
        \[ F_{B/A} = F_{A/B} = K\,\frac{q_A\,q_B}{d^2} = K\,\frac{(4q)(q)}{d^2}. \]
        Elles ont donc \textbf{la même intensité} (principe de l'action et de la réaction) : la
        charge $4$ fois plus grande de $A$ n'y change rien, car c'est le \emph{produit} $q_A q_B$
        qui compte.
  \item Les deux boules ont même masse et subissent une force de même intensité : elles s'écartent
        du \textbf{même angle}, $\alpha_A = \alpha_B$. \emph{Non}, la boule la plus chargée ne
        s'écarte pas davantage — c'est le piège classique de cette situation.
\end{enumerate}
\end{corrige}

\end{document}
